% page 315 du fac-simile (image p-314 du scan)
% bloc 165.1 x 242.0 mm ; marge gauche 36.9 mm
\pgc{15.240mm}{0.000mm}{
\l{\cel{53}307}
\l{}
\l{\sou{koriono.}\cel{2}(\sou{anat.)}\cel{3}Envelopilo extera di la feto. - DEFIS.}
\l{}
\l{\sou{korizo.}\cel{2}(\sou{biol.})\cel{2}Duimesko o duoplesko di ula parti, per for-}
\l{\cel{3}maco di organi supernombra. -}
\l{}
\l{\sou{korko.}\cel{2}Ta strato sponjatra e lejera di speco di querko (L.}
\l{\cel{3}\sou{quercum suber}), de qua onu facas botelo-stopili, flotaceri,}
\l{\cel{3}e c. - DE.}
\l{}
\l{\sou{kormorano.}\cel{2}(\sou{zool.)}\cel{2}Ucelo qua formacas grupo ek la klaso}
\l{\cel{3}\textquotedbl{}palmipedi\textquotedbl{}, qua vivas rive di la aquo-loki (lagi, mari,}
\l{\cel{3}riveri) e qua su nutras ye fishi. - L.\cel{1}\sou{phalacrocorax (carbo)}.}
\l{\cel{3}- DEFI.}
\l{}
\l{\sou{korno.}\cel{2}I. Surkreskajo konatra, rekta o kurva, qua kreskas sur}
\l{\cel{3}la fronto di la rumineri, ordinare che la maskulo, e formacas}
\l{\cel{3}quaza prolonguro di la osto frontala. - II. (\sou{analogie)}.}
\l{\cel{3}Surkreskajo, kornatra,qua existas sur la kapo di ula animali.}
\l{\cel{3}Anke : pedunkli retrotirebla, qua existas sur la kapo di la}
\l{\cel{3}limaki e qua portas la organi vidala. - III. To quo finas}
\l{\cel{3}per punto simila a korno. - IV. La materio tenaca,lameloza,}
\l{\cel{3}de qua facesas la korni di la rumineri. - V. Korno kavigita e}
\l{\cel{3}truizita, quan uzis la muton-pastori, la muton-gardisti, la}
\l{\cel{3}chaseri, e c. por signalar, advokar. - VI. (\sou{muziko)}. Instru-}
\l{\cel{3}mento suflala, ek kupro, tordita adspirale, di qua un extre-}
\l{\cel{3}majo recevas boko-peco (mikra kupeto ek metalo, ek ivoro,}
\l{\cel{3}e c.) e di qua la altra extremajo formacas kono expansita quan}
\l{\cel{3}onu nomizis \textquotedbl{}funelo\textquotedbl{}. - EFIS.}
\l{}
\l{\sou{kornajar.(netrans}.) (\sou{aludante kavalo, asno)}. Produktar ta ster-}
\l{\cel{3}toro quan kreas morbo di la organi respirala. - DEFIRS.}
\l{}
\l{\sou{kornalino}. (\sou{mineral.)}\cel{3}Agato mi-diafana, plu o min obskure-reda,}
\l{\cel{3}de qua onu facas siglili, ringi fingrala, e c. - DEFIS.}
\l{}
\l{\sou{kornamuzo.}\cel{2}(\sou{muziko).}\cel{2}Muzik-instrumento rustika, di qua la son-}
\l{\cel{3}koloro esas nazo-sonanta, facita ek felto-sako, aero-tanko,}
\l{\cel{3}quan onu inflas per la boko, mediace tubo, ed ek tri tubi}
\l{\cel{3}muntita sur la felto-sako, de qui du emisas noto kontinua,}
\l{\cel{3}e de qui la triesma esas provizita ye boko-peco langetoza}
\l{\cel{3}aden qua suflas la pleanto, e borita ye trui sur sui lu}
\l{\cel{3}promenigas sua fingri por produktar diversa noti. - FIS.}
\l{}
\l{\sou{korneo}. (\sou{anat}.)\cel{2}Membrano qua tapisas la dopa parto di la}
\l{\cel{3}okulo. DEFIS.}
\l{}
\l{\sou{kornelo. (bot.}) Frukto de arboro di qua la ligno esas tre harda,}
\l{\cel{3}qua formacas la speco tipa di la familio \textquotedbl{}kornacei\textquotedbl{}, e qua}
\l{\cel{3}similesas mikra pruno reda. - L.\cel{1}\sou{cornus mas.}}
\l{}
\l{korneto.\cel{2}I. Mikra korno, rustika, korno ek terakoto, garnisita}
\l{\cel{3}ye boko-peco, per qua la pueri su amuzas, pleante lu lor}
\l{\cel{3}karnavalo. - Granda fluto quan onu uzas en la kori por sub-}
\l{\cel{3}tenar la voci. - II.To quo esas kornatra. - DEFIS.}
}
